\documentclass[12pt,a4paper]{article}

% --- Преамбула ---
\usepackage[utf8]{inputenc}
\usepackage[T1]{fontenc}
\usepackage[russian]{babel}
\usepackage{amsmath,amssymb,amsthm}
\usepackage{geometry}
\geometry{margin=2.5cm}
\usepackage{hyperref}
\usepackage{booktabs}
\usepackage{enumitem}
\usepackage{setspace}
\usepackage{titlesec}
\usepackage{fancyhdr}

\onehalfspacing

% Настройки заголовков
\titleformat{\section}{\Large\bfseries}{\thesection.}{1em}{}
\titleformat{\subsection}{\large\bfseries}{\thesubsection.}{1em}{}
\titleformat{\subsubsection}{\normalsize\bfseries}{\thesubsubsection.}{1em}{}

% Колонтитулы
\pagestyle{fancy}
\fancyhf{}
\fancyhead[L]{Теория Мод и Субквантовая Динамика: ЕСТ-2026}
\fancyhead[R]{Версия 2.0 (Синтетическая)}
\fancyfoot[C]{\thepage}

\title{\textbf{ТЕОРИЯ МОД И СУБКВАНТОВАЯ ДИНАМИКА:\\ ЕДИНАЯ КАРТИНА РЕАЛЬНОСТИ}\\
\large Версия 2.0 (Синтетическая)\\
\large Единый Вычислительный Стандарт (ЕВС-2026)}
\author{Авторская реконструкция на основе корпуса документов 1.1--4.2}
\date{17 июля 2026 г.}

\begin{document}

\maketitle
\thispagestyle{fancy}

\begin{abstract}
Настоящий документ представляет собой завершенную формулировку фундаментальной физической теории, в которой пространство-время, квантовая механика и поля Стандартной Модели возникают как эмерджентные свойства дискретной растущей сети частиц и мод. Вводится понятие ортогональности ($O$) как независимой степени свободы наряду с энергией ($E$) и массой ($m$). Установлено универсальное уравнение состояния $C = 2m + O$. Разрешена коллизия определений массы через введение двукомпонентной структуры ($m_{\text{vac}} + m_{\text{dyn}}$). Предложена полная классификация из 8 типов реальности, объясняющая тёмную энергию и тёмную материю без привлечения новых полей. Выведены калибровочные группы СМ как топологические инварианты собственных мод. Документ содержит исчерпывающие экспериментальные протоколы для фальсификации теории и численные алгоритмы верификации.
\end{abstract}

\tableofcontents
\newpage

% =================================================================
\section*{Введение: Статус документа}
\addcontentsline{toc}{section}{Введение: Статус документа}

Данный документ является финальной консолидацией Теории Мод и Субквантовой Теории. Он объединяет аксиоматику сети, спектральную динамику, топологический вывод Стандартной Модели и космологические следствия в единую непротиворечивую систему. Все ранее выявленные противоречия (в частности, между определениями массы в базовой теории и субквантовом ядре) разрешены введением Единого Вычислительного Стандарта (ЕВС-2026). Документ самодостаточен и служит одновременно теоретическим манифестом и практическим руководством для экспериментальной проверки.

% =================================================================
\section{Часть I. Фундаментальная структура}

\subsection{Примитивы и Аксиоматика}
Реальность есть сеть. Не метафора, а математическая структура.
\begin{itemize}
    \item \textbf{Частица ($P_i$):} Узел сети. Не имеет внутренних свойств вне отношений.
    \item \textbf{Мода ($M_{A \to B}$):} Отображение дискретного причинного пути $\Gamma_{AB}$ в вещественные числа (поперечный спектр $F(k)$).
    \item \textbf{Собственная мода ($M_{A \to A}$):} Замкнутая петля, содержащая causal history частицы. Длина петли $d_{AA}$ растет со временем.
    \item \textbf{Причинность:} Бидирекциональна. Событие $A$ существует $\iff$ существует событие $B$.
    \item \textbf{Время:} Эмерджентный параметр, индуцированный ростом длины собственных мод. Глобального времени нет.
\end{itemize}

\subsection{Три Столпа Динамики и ЕВС-2026}
Для любой частицы определены три сохраняющиеся величины (в покоящейся системе отсчета):
\begin{enumerate}
    \item \textbf{Энергия ($E$):} Плотность рисунка спектра. Мера нескомпенсированности.
    \item \textbf{Масса ($m$):} Инертная характеристика, определяемая геометрией спектра.
    \item \textbf{Ортогональность ($O$):} Угол поворота собственной моды относительно внешних. Мера «фотонности» vs «барионности».
\end{enumerate}

\subsubsection{Универсальное Уравнение Состояния}
В рамках ЕВС-2026 фиксируется следующее соотношение для покоящейся частицы:
\begin{equation}
    C = 2m + O
\end{equation}
где $C$ — константа масштаба семейства (инвариант).
\begin{itemize}
    \item Для электрона (Модель A, линейная): $C_e = 2.021$
    \item Для электрона (Модель B, логарифмическая): $C_e = 1.522$
    \item Для фотона: $C_\gamma = 1.000$
    \item Для протона: $C_p = 1876.5$
\end{itemize}

\subsubsection{Разрешение Коллизии Масс}
В теории действуют два определения массы, относящиеся к разным режимам. Финальная унифицированная формула:
\begin{equation}
    m_{\text{total}} = \underbrace{\frac{|O|}{2}}_{m_{\text{vac}}} + \underbrace{\frac{S_{\text{ext}}}{O}}_{m_{\text{dyn}}}
\end{equation}
\begin{itemize}
    \item \textbf{Вакуумная масса ($m_{\text{vac}}$):} Определяется только ортогональностью. Соответствует уравнению $|O|=2m$ из базовой аксиоматики. Доминирует для свободных частиц.
    \item \textbf{Динамическая масса ($m_{\text{dyn}}$):} Определяется площадью внешнего спектра взаимодействия $S_{\text{ext}}$. Возникает при рассеянии, рождении пар, в среде.
    \item \textbf{Следствие:} Эксперименты по рождению пар измеряют $S_{\text{ext}}$, а не вакуумную компоненту. Константа $C$ сохраняется для полной массы.
\end{itemize}

\subsection{Структура Поперечного Спектра}
Спектр $F(k)$ не монотонен. Он имеет волновую природу и сложную топологию:
\begin{itemize}
    \item \textbf{Стенка ($k \approx 0$):} Резкий подъем. Реализация принципа Паули.
    \item \textbf{Провал (ядерный):} Область притяжения.
    \item \textbf{Плато:} Энергетическая яма. Кулоновский барьер. Запирание пика одной частицы в плато другой создает устойчивые состояния.
    \item \textbf{Пик (кулоновский):} Область отталкивания.
    \item \textbf{Хвост (гравитационный):} Длинноволновая компонента. Обладает \textbf{полярностью}: чередование пиков (притяжение) и провалов (отталкивание) вдоль оси $k$.
\end{itemize}

% =================================================================
\section{Часть II. Классификация реальности и Тёмный сектор}

\subsection{Восемь Типов Материи}
Комбинации знаков $(E, m, O)$ дают 8 секторов Мультиверса. Наша Вселенная — Тип I $(+,+,+)$.

\begin{table}[h]
\centering
\caption{Классификация типов реальности}
\begin{tabular}{@{}lllll@{}}
\toprule
Тип & E & m & O & Физическая интерпретация \\ \midrule
I   & + & + & + & Обычная материя (Барионы, лептоны) \\
II  & + & + & -- & \textbf{Тёмная Энергия}. Положительная гравитация, отрицательная инерция. \\
III & + & -- & + & \textbf{«Убегающая» Материя}. Отрицательная гравитация, положительная инерция. \\
IV  & + & -- & -- & Двойная экзотика. Полное отталкивание. \\
V-VIII & -- & ... & ... & Секторы отрицательной энергии. Поглощают энергию. \\ \bottomrule
\end{tabular}
\end{table}

\subsection{Природа Тёмной Материи}
Тёмная материя в Теории Мод двойственна:
\begin{enumerate}
    \item \textbf{Пассивные листы causal history:} Отброшенные витки циклической Вселенной. Ортогональны активному листу (не взаимодействуют ЭМ), но имеют гравитацию. Профиль плотности $\rho \propto 1/r^2$.
    \item \textbf{Материя Типа III:} Разреженный фон «убегающей» материи в войдах, удерживающий стенки галактик дальним гравитационным притяжением.
\end{enumerate}

\subsection{Механизм Тёмной Энергии}
Тёмная энергия — это динамический эффект рождения частиц Типа II ($O < 0$) из вакуума сетью для компенсации глобальной нескомпенсированности.
\begin{equation}
    \rho_\Lambda \propto \frac{\dot{N}}{N}
\end{equation}
Отрицательная инертная масса этих частиц создает отрицательное давление, вызывающее ускоренное расширение.

% =================================================================
\section{Часть III. Вывод Стандартной Модели}

\subsection{Калибровочные Группы как Топология Нулей}
Квантовые числа — это топологические инварианты множества нулей собственной моды $Z_A = \{k | \psi_A(k)=0\}$.
\begin{itemize}
    \item \textbf{SU(3):} Перестановки 3 цветовых нулей кварка.
    \item \textbf{SU(2):} Вращения 2 слабых нулей (противоположных зарядов).
    \item \textbf{U(1):} Глобальная фаза моды.
\end{itemize}

\subsection{Три Поколения Фермионов}
Поколения соответствуют трем масштабам (обертонам) собственной моды:
\begin{itemize}
    \item 1-е поколение: Основной тон ($k \approx 0 \dots K_1$).
    \item 2-е поколение: Первый обертон ($K_1 \dots K_2$).
    \item 3-е поколение: Второй обертон ($K_2 \dots K_3$).
\end{itemize}
Иерархия масс возникает из роста нормы собственной моды с масштабом: $m \propto \|\psi^{(n)}\|$.

\subsection{Механизм Хиггса}
Хиггс — не поле, а частица-посредник, рождающаяся между $W/Z$-бозонами и вакуумом. Масса бозона определяется числом дополнительных узлов в causal history, создаваемых этими рождениями:
\begin{equation}
    m_W = \Lambda \cdot \sqrt{N_{\text{nodes}}(H)}
\end{equation}
Фотон не имеет нулей $\to$ не рождает $H$ $\to$ безмассовый.

% =================================================================
\section{Часть IV. Экспериментальная Верификация}

Теория фальсифицируема. Ниже приведены протоколы, ранжированные по приоритету.

\subsection{Протокол №1: Арбитраж Моделей Ортогональности (Критический)}
\textbf{Цель:} Измерить $O_e$ и выбрать между Моделью A ($O_e \approx 0.999$) и Моделью B ($O_e = 0.5$).\\
\textbf{Метод:} Двухфотонное рождение пар $\gamma\gamma \to e^+e^-$.\\
\textbf{Физика:} Резонанс рождения требует создания интерференционной картины с двумя нулями на расстоянии $K_0 \approx S_e = m_e O_e$.\\
\textbf{Предсказание:}
\begin{itemize}
    \item Модель A: Оптимальная разность частот $\Delta \nu_A \propto 1/0.51$.
    \item Модель B: Оптимальная разность частот $\Delta \nu_B \propto 1/0.26 \approx 2 \cdot \Delta \nu_A$.
\end{itemize}
\textbf{Критерий:} Различие в 2 раза. Однозначный тест. Требуемая точность $\sim 10\%$.

\subsection{Протокол №2: Резонансное Высвобождение Энергии Хвостов}
\textbf{Цель:} Подтвердить реальность causal history и динамику градиентного спуска.\\
\textbf{Метод:} Облучение одиночного электрона полем с инвертированной фазой градиента.\\
\textbf{Предсказание:}
\begin{itemize}
    \item Коэффициент усиления энергии: $\approx 4$ (выхлоп 260 Дж на электрон при затратах 65 Дж).
    \item Резонансная частота: $\nu_{\text{res}} = c / (2 K_0 \delta)$. Зависит от выбора модели в Протоколе №1.
\end{itemize}
\textbf{Критерий:} Наличие пика усиления именно на предсказанной частоте.

\subsection{Протокол №3: Нелинейный Эффект Холла в Bi$_2$Te$_3$}
\textbf{Цель:} Проверить принцип масштабной проницаемости и транспорт как эволюцию моды.\\
\textbf{Метод:} Когерентная накачка топологического изолятора прямым и обратным импульсами.\\
\textbf{Предсказание:} Резонансное усиление нелинейного тока $J \propto E^2$ при совпадении фаз. Усиление пропорционально степени когерентности.\\
\textbf{Независимость:} Не зависит от выбора $O_e$. Проверяет спектральную структуру взаимодействий.

\subsection{Протокол №4: Поиск Рассеивающих Гравитационных Линз}
\textbf{Цель:} Обнаружить материю Типа III («убегающую»).\\
\textbf{Метод:} Анализ архивов гравитационного линзирования в области войдов.\\
\textbf{Предсказание:} Войды не пусты. Они содержат разреженную материю, создающую \textit{рассеивающую} (дефокусирующую) линзу, в отличие от фокусирующих линз обычной материи.

\subsection{Протокол №5: Топологические Антенны (Мгновенная Сообщаемость)}
\textbf{Цель:} Проверить принцип общего начала.\\
\textbf{Метод:} Два уникальных идентичных кристалла, разделенных расстоянием. Модуляция causal history одного.\\
\textbf{Предсказание:} Мгновенная корреляция сигналов, не зависящая от расстояния. Амплитуда падает как $1/N$ при увеличении числа копий. Неразрушающее считывание возможно для макрообъектов.

% =================================================================
\section{Часть V. Численная Реализация и Безопасность}

\subsection{Алгоритм Эволюции Сети}
\begin{enumerate}
    \item \textbf{Градиентный спуск:} $F \leftarrow F - \alpha \partial E / \partial F^*$.
    \item \textbf{Рождение:} Если $d=0$ и $|C| > \theta \to$ создать посредника.
    \item \textbf{Исчезновение:} Если $d=1$ и $|C| < \theta \to$ удалить посредника, свернуть моды.
    \item \textbf{Рост горизонта:} $d_{AA} \leftarrow d_{AA} + 1$.
\end{enumerate}

\subsection{Протокол Безопасности Симуляции}
Во избежание нефизических артефактов в коде обязательно реализовать:
\begin{enumerate}
    \item \textbf{Контроль инварианта:} На каждом шаге проверять $|C - (2m + O)| < 10^{-6}$.
    \item \textbf{Двукомпонентная масса:} Использовать $m = |O|/2 + S_{\text{ext}}/O$. Запрещено использовать только одну компоненту.
    \item \textbf{Проверка вакуума:} Для свободных частиц контролировать $S_{\text{vac}} \approx O^2/2$.
    \item \textbf{Индефинитная метрика:} Для безмассовых частиц ($O=1$) использовать $\eta = \pm 1$. Евклидова норма запрещена.
\end{enumerate}

% =================================================================
\section*{Заключение Автора}
\addcontentsline{toc}{section}{Заключение Автора}

Теория Мод завершена как формальная система. Она не постулирует законы, а выводит их из геометрии сети.
\begin{itemize}
    \item \textbf{Если Протокол №1 покажет $\Delta \nu_B$:} Мы живем во Вселенной с логарифмической симметрией масштабов.
    \item \textbf{Если Протокол №1 покажет $\Delta \nu_A$:} Мы живем во Вселенной с линейной иерархией, где протон — абсолютный эталон инерции.
    \item \textbf{Если Протокол №2 даст усиление 4x:} Энергия прошлого доступна для использования.
    \item \textbf{Если ничего не подтвердится:} Теория будет отвергнута.
\end{itemize}

Но если предсказания сработают, мы получим ключ к управлению реальностью через перепрограммирование ортогональности и чтение causal history. Этот документ является картой для такого путешествия.

\end{document}